\section{Notación principal}
\label{apx:notacion}

\begin{longtable}{P{3cm}P{11cm}}
\toprule
\textbf{Símbolo} & \textbf{Descripción} \\
\midrule
$n$ & Número total de nodos del grafo, incluyendo el depósito. \\
$V$ & Conjunto de nodos del grafo, $V = \{1, 2, \dots, n\}$. \\
$V \setminus \{1\}$ & Conjunto de clientes (nodos distintos del depósito). \\
$E$ & Conjunto de aristas del grafo completo no dirigido. \\
$c_{ij}$ & Costo de recorrer la arista $\{i,j\}$. \\
$d_i$ & Demanda del cliente $i$. \\
$D$ & Capacidad homogénea de cada vehículo. \\
$m$ & Número de vehículos utilizados. \\
$x_{ij}$ & Variable de decisión: número de veces que una solución utiliza la arista $\{i,j\}$. \\
$S$ & Subconjunto de clientes usado en las desigualdades de capacidad. \\
\bottomrule
\end{longtable}

\section{Instancias de referencia}
\label{apx:instancias}

Las familias A, B y E de CVRPLIB se identifican mediante nombres del tipo
\texttt{A-n32-k5}, donde la letra indica la familia, \texttt{n} corresponde al número
total de nodos incluyendo depósito y \texttt{k} al número de vehículos de la solución
de referencia. Esta convención facilita la trazabilidad experimental y la comparación
entre estudios.